\subsection{误差分析：连续留段与经验覆盖}
\label{sec:error-analysis}

碳化硅的谱形预测有所改善，区间覆盖却出现缺口。误差检验先拆开这两件事，再用材料分项和合成回收辨明测厚误差的参照。

原型比较曾给出主方法误差 $1.1717$；改用相同连续留段复核后，正式误差为 $1.2074$，低于二次趋势的 $1.4404$，改善 $16.17\%$，故采用正式评分衡量预测增益。上述误差均为无量纲的标准化反射率均方根误差，图\ref{fig:q2-residual}中的局部残差对应这一比较。
% src:交接/实验记录.json:528.尝试;528.现象;528.决定
% src:求解/问题2/结果/验证.json:三路线原型对比.双角色散变投影;正式主方法分数_连续留段标准化均方根误差;正式二次趋势基线分数_连续留段标准化均方根误差;正式主方法相对正式二次趋势改善_百分比

预测增益没有消除漏盖：测试反射率仅有 $157/320$ 点落入预测区间，经验覆盖率为 $49.06\%$。它衡量留出观测落入区间的比例；图\ref{fig:q2-coverage}与表\ref{tab:09}将这一比例和区间宽度并列。问题二采用校准残差的 $75\%$ 分位确定半宽，再由留出的测试块统计实际覆盖。
% src:求解/问题2/结果/区间覆盖.json:测试经验覆盖率;测试点数;校准经验分位点;构造方法
% src:求解/问题2/结果/验证.json:分块

\begingroup
\small
\setlength{\LTpre}{0.35em}
\setlength{\LTpost}{0.35em}
\begin{longtable}{>{\centering\arraybackslash}p{0.15\textwidth}>{\centering\arraybackslash}p{0.14\textwidth}>{\centering\arraybackslash}p{0.14\textwidth}>{\centering\arraybackslash}p{0.15\textwidth}>{\centering\arraybackslash}p{0.12\textwidth}>{\centering\arraybackslash}p{0.10\textwidth}}
\caption{正式实测反射率的点覆盖与整块覆盖}\label{tab:09}\\
\toprule
对象 & 分位/名义 & 点覆盖 & 全宽均值 & 整块覆盖 & 漏盖点\\
\midrule
\endfirsthead
\toprule
对象 & 分位/名义 & 点覆盖 & 全宽均值 & 整块覆盖 & 漏盖点\\
\midrule
\endhead
碳化硅 & $75\%$/--- & $157/320$ & $0.7424$ & $3/8$ & $163$\\
两材料 & $90\%/90\%$ & $473/640$ & $3.0875$ & $9/16$ & $167$\\
\bottomrule
\end{longtable}
% src:求解/问题2/结果/区间覆盖.json:校准经验分位点;构造方法;测试经验覆盖率;测试点数
% src:求解/问题2/结果/验证.json:分块
% src:求解/问题3/结果/回炉轮3候选/仲裁闭环/20260910_131501_81623_建模公式/厚度结果.json:预测区间.4.经验覆盖率;预测区间.4.平均区间宽度_比例;预测区间.5.经验覆盖率;预测区间.5.平均区间宽度_比例;预测区间.6.经验覆盖率;预测区间.6.平均区间宽度_比例;预测区间.7.经验覆盖率;预测区间.7.平均区间宽度_比例;预测区间.20.经验覆盖率;预测区间.20.平均区间宽度_比例;预测区间.21.经验覆盖率;预测区间.21.平均区间宽度_比例;预测区间.22.经验覆盖率;预测区间.22.平均区间宽度_比例;预测区间.23.经验覆盖率;预测区间.23.平均区间宽度_比例;预测区间.36.经验覆盖率;预测区间.36.平均区间宽度_比例;预测区间.37.经验覆盖率;预测区间.37.平均区间宽度_比例;预测区间.38.经验覆盖率;预测区间.38.平均区间宽度_比例;预测区间.39.经验覆盖率;预测区间.39.平均区间宽度_比例;预测区间.52.经验覆盖率;预测区间.52.平均区间宽度_比例;预测区间.53.经验覆盖率;预测区间.53.平均区间宽度_比例;预测区间.54.经验覆盖率;预测区间.54.平均区间宽度_比例;预测区间.55.经验覆盖率;预测区间.55.平均区间宽度_比例;预测区间.4.名义覆盖率
% src:求解/问题3/结果/回炉轮3候选/仲裁闭环/20260910_131501_81623_建模公式/厚度结果.json:逐块评分.4.测试点数;逐块评分.5.测试点数;逐块评分.6.测试点数;逐块评分.7.测试点数;逐块评分.20.测试点数;逐块评分.21.测试点数;逐块评分.22.测试点数;逐块评分.23.测试点数;逐块评分.36.测试点数;逐块评分.37.测试点数;逐块评分.38.测试点数;逐块评分.39.测试点数;逐块评分.52.测试点数;逐块评分.53.测试点数;逐块评分.54.测试点数;逐块评分.55.测试点数
% src:求解/问题3/结果/回炉轮3候选/仲裁闭环/20260910_131501_81623_建模公式/厚度结果.json:预测区间.4.经验覆盖率;预测区间.4.平均区间宽度_比例;预测区间.5.经验覆盖率;预测区间.5.平均区间宽度_比例;预测区间.6.经验覆盖率;预测区间.6.平均区间宽度_比例;预测区间.7.经验覆盖率;预测区间.7.平均区间宽度_比例;预测区间.20.经验覆盖率;预测区间.20.平均区间宽度_比例;预测区间.21.经验覆盖率;预测区间.21.平均区间宽度_比例;预测区间.22.经验覆盖率;预测区间.22.平均区间宽度_比例;预测区间.23.经验覆盖率;预测区间.23.平均区间宽度_比例;预测区间.36.经验覆盖率;预测区间.36.平均区间宽度_比例;预测区间.37.经验覆盖率;预测区间.37.平均区间宽度_比例;预测区间.38.经验覆盖率;预测区间.38.平均区间宽度_比例;预测区间.39.经验覆盖率;预测区间.39.平均区间宽度_比例;预测区间.52.经验覆盖率;预测区间.52.平均区间宽度_比例;预测区间.53.经验覆盖率;预测区间.53.平均区间宽度_比例;预测区间.54.经验覆盖率;预测区间.54.平均区间宽度_比例;预测区间.55.经验覆盖率;预测区间.55.平均区间宽度_比例;预测区间.4.名义覆盖率
% src:求解/问题3/结果/回炉轮3候选/仲裁闭环/20260910_131501_81623_建模公式/厚度结果.json:逐块评分.4.测试点数;逐块评分.5.测试点数;逐块评分.6.测试点数;逐块评分.7.测试点数;逐块评分.20.测试点数;逐块评分.21.测试点数;逐块评分.22.测试点数;逐块评分.23.测试点数;逐块评分.36.测试点数;逐块评分.37.测试点数;逐块评分.38.测试点数;逐块评分.39.测试点数;逐块评分.52.测试点数;逐块评分.53.测试点数;逐块评分.54.测试点数;逐块评分.55.测试点数
\endgroup

表中全宽以反射率百分点计，两材料行合计硅与碳化硅的完整往返场预测。整块覆盖要求该块全部观测落入同一点预测带，漏盖点保留逐点预测失败的次数。此处统计的整块覆盖不是另构造的块同时预测带覆盖；连续谱点彼此相关，点覆盖与整块覆盖分别检验该点预测带在单个观测及整段观测上的表现。

材料分项揭示了总体均值掩盖的冲突。完整往返场误差 $1.0231$ 高于两束的 $0.9812$；硅附件三、四各有 $2$ 个负改善块。表\ref{tab:supp-q3-blocks}逐项保留差值，图\ref{fig:16}与图\ref{fig:17}分别展示局部残差和块间方向；局部优势没有转化为总体预测改善。
% src:求解/问题3/结果/回炉轮3候选/仲裁闭环/20260910_131501_81623_建模公式/厚度结果.json:分材料验证.硅.逐角逐折逐块.0.完整相对两束下降_%;分材料验证.硅.逐角逐折逐块.1.完整相对两束下降_%;分材料验证.硅.逐角逐折逐块.2.完整相对两束下降_%;分材料验证.硅.逐角逐折逐块.3.完整相对两束下降_%;分材料验证.硅.逐角逐折逐块.4.完整相对两束下降_%;分材料验证.硅.逐角逐折逐块.5.完整相对两束下降_%;分材料验证.硅.逐角逐折逐块.6.完整相对两束下降_%;分材料验证.硅.逐角逐折逐块.7.完整相对两束下降_%
% src:求解/问题3/结果/回炉轮3候选/仲裁闭环/20260910_131501_81623_建模公式/厚度结果.json:分材料验证.硅.逐折汇总.0.平均标准化均方根误差.两束;分材料验证.硅.逐折汇总.0.平均标准化均方根误差.完整往返;分材料验证.硅.逐折汇总.0.平均标准化均方根误差.训练均值;分材料验证.硅.逐折汇总.0.平均标准化均方根误差.二次趋势;分材料验证.硅.逐折汇总.1.平均标准化均方根误差.两束;分材料验证.硅.逐折汇总.1.平均标准化均方根误差.完整往返;分材料验证.硅.逐折汇总.1.平均标准化均方根误差.训练均值;分材料验证.硅.逐折汇总.1.平均标准化均方根误差.二次趋势;分材料验证.碳化硅.逐折汇总.0.平均标准化均方根误差.两束;分材料验证.碳化硅.逐折汇总.0.平均标准化均方根误差.完整往返;分材料验证.碳化硅.逐折汇总.0.平均标准化均方根误差.训练均值;分材料验证.碳化硅.逐折汇总.0.平均标准化均方根误差.二次趋势;分材料验证.碳化硅.逐折汇总.1.平均标准化均方根误差.两束;分材料验证.碳化硅.逐折汇总.1.平均标准化均方根误差.完整往返;分材料验证.碳化硅.逐折汇总.1.平均标准化均方根误差.训练均值;分材料验证.碳化硅.逐折汇总.1.平均标准化均方根误差.二次趋势

第三问原第三版往返场估计器在低损耗、强损耗和无高阶负对照的$9$例匹配情景中，最大厚度相对误差绝对值为$0.0521\%$，但近优条件包络仅有$1/9$包含真厚度。表\ref{tab:q3-known-thickness}保留全部已知厚度案例；这组固定光学参数检验衡量厚度搜索的回收能力，与现行候选整合后的实测结果分列。
% src:求解/问题3/结果/公平对照.json:已知厚度验证.完成匹配情景数;已知厚度验证.案例.0.相对误差_%;已知厚度验证.案例.1.相对误差_%;已知厚度验证.案例.2.相对误差_%;已知厚度验证.案例.3.相对误差_%;已知厚度验证.案例.4.相对误差_%;已知厚度验证.案例.5.相对误差_%;已知厚度验证.案例.6.相对误差_%;已知厚度验证.案例.7.相对误差_%;已知厚度验证.案例.8.相对误差_%;已知厚度验证.条件包络经验包含率

色散失配时，厚度相对误差绝对值分别为$3.71\%$、$8.96\%$和$7.35\%$。这组检验将估计模型中的色散设为零，偏差随之放大，说明匹配光学参数时的小回收误差不能代表光学关系失配时的测厚表现。
% src:求解/问题3/结果/公平对照.json:已知厚度验证.案例.9.相对误差_%;已知厚度验证.案例.10.相对误差_%;已知厚度验证.案例.11.相对误差_%;已知厚度验证.案例.9.口径

合成近优包络收集固定生成条件下训练目标接近最优的厚度解；实测多情景包络还联合区段切分和光学条件变化，构造对象不同。前者的真值包含比例是这组模拟的经验计数，后者没有实测厚度真值可供同样计数，二者均未建立概率覆盖保证。

模型失配也放大了问题一的厚度误差：平均绝对相对误差由 $0.0008\%$ 升至 $0.3028\%$（图\ref{fig:08}），共同匹配组的常折射率峰距基线失败 $0$ 次。该项检验与第三问的直接厚度回收分别报告。合成回收以已知厚度为参照核对求解过程，光学参数扰动则考察固定观测下的厚度变化，下节沿这一方向展开重估比较。
% src:求解/问题1/结果/合成验证.json:主指标.数值;分组汇总.模型失配.平均绝对相对误差_百分比;常折射率峰距基线.失败数
